%%%%%%%%%%  scan idx 356  |  folio 343  %%%%%%%%%%

Donc \(\lambda^p(N+L,x)\) est le coefficient de \(t^p\) dans la série :
\[
  \frac{\displaystyle\sum_{k\ge1}\lambda^k(N,x)
        (1+\cdots+L^{k-1})t^k}
       {\displaystyle 1+\sum_{k\ge1}\lambda^k(N,x)
        \lambda_{-1}(N)L^k t^k}.
\]

\textbf{Proposition 5.7.} Soient \(K\) un anneau binomial, \(A\) une \(K\)-\(\lambda\)-algèbre augmentée, \(N\) un élément de \(A\) tel que \(\lambda^k(N)=0\) pour \(k>d\). Alors pour tout \(x\in A\) et tout \(p\ge1\), on a :
\begin{equation}\tag{5.7.1}
  \gamma^p(N,x)\in\Fil^{p-d}A.
\end{equation}

Soit \(A_0\) la \(K\)-\(\lambda\)-algèbre engendrée par des générateurs \(N'_0,X'_0\), soumis aux relations \(\gamma^k(N'_0)=0\) pour \(k>d\). On définit l'homomorphisme d'augmentation \(A_0\xrightarrow{\epsilon}K\) par \(\epsilon(N'_0)=\epsilon(X'_0)=0\). On sait que \(A_0\) est l'anneau des polynômes à coefficients dans \(K\) par rapport aux indéterminées \(\gamma^i(N'_0)\), avec \(1\le i\le d\), et \(\gamma^j(X'_0)\), avec \(j\ge1\), et que la \(\lambda\)-filtration sur \(A_0\) est identique à la filtration par le poids, en affectant le poids \(i\) à \(\gamma^i(N'_0)\), \(\gamma^i(X'_0)\) (4.12 et 4.15).

Posons \(N_0=N'_0+d\), \(X_0=X'_0+\epsilon(x)\). Nous allons d'abord montrer 5.7 pour \(N_0\) et \(X_0\). D'après ce qui précède, (5.7.1) est équivalent à :
\[
  \gamma^p(N_0,X_0)\cdot\gamma^d(N'_0)\in\Fil^pA_0.
\]
Or \(\gamma^d(N'_0)=(-1)^d\lambda_{-1}(N_0)\)~; cette relation s'écrit donc :
\[
  (-1)^d\gamma^p(X_0\gamma^d(N_0))\in\Fil^pA_0
\]
d'après (5.5.4). Cette relation est bien vérifiée, puisque \(\gamma^d(N_0)\) est d'augmentation nulle.

D'après 4.7, il existe un \(K\)-\(\lambda\)-homomorphisme \(A_0\longrightarrow A\) compatible aux augmentations envoyant \(N'_0\) sur \(N-d\), et \(X'_0\) sur \(x-\epsilon(x)\), donc \(N_0\) sur \(N\) et \(X_0\) 
